\subsection{赋值的高}

设 G 是全序群。对 G 的两个孤立子群 H 与 H'，二者之一包含于另一个：否则会存在 H 的不属于 H' 的正元素 x 以及 H' 的不属于 H 的正元素 \(x'\)；例如设 \(x \geqslant x'\) ；由于 H 是孤立的，\(x' \in H\) ，这就得到矛盾。

这一点也可以由第 3 节命题 4 与第 1 节命题 1 的推论得出，只需注意到每个全序群都是某个赋值的阶群（§ 3 第 4 节例 6）。

定义 2. 设 G 是全序群。若 G 的异于 G 的孤立子群的个数有限且等于 n，则称 G 的高为 n。~若这个数是无限的，则称 G 的高为无限。

\textbf{例}\par

\begin{enumerate}
\def\labelenumi{(\arabic{enumi})}
\item
  群 \(\mathbf{G} = \{0\}\) 的高为 0。
\item
  群 \(\mathbf{Z}\) 与 \(\mathbf{R}\) 的高为 1。
\item
  设 G 是全序群，H 是 G 的孤立子群。若用 \(h(\mathrm{H})\) 与 \(h(\mathrm{G}/\mathrm{H})\) 分别表示全序群 H 与 G/H 的高，则
\end{enumerate}

\[
h (\mathbf {G}) = h (\mathbf {H}) + h (\mathbf {G} / \mathbf {H}),\tag{2}
\]

因为 G 的孤立子群的集合按包含关系是全序的。特别地，若 G 是两个全序群 H 与 H' 的字典积，则

\[
h (\mathbf {G}) = h (\mathbf {H}) + h \left(\mathbf {H} ^ {\prime}\right)\tag{3}
\]

（参见~第 2 节例 2）；于是字典积 \(\mathbf{Z} \times \mathbf{Z}\) 的高为 2。

另一方面，\(\mathbf{Z} \times \mathbf{Z}\) 按嵌入 \(\mathbf{R}\) 的序排序时（参见~§ 3 第 4 节例 6 末尾），其高等于 1（参见~下面的命题 8）。

定义 3. 赋值的阶群的高称为该赋值的高。

例如，离散赋值的高为 1。只有非真赋值的高为 0。命题 1 与命题 4 蕴含：

命题 5. 赋值的高等于其环中非零素理想的个数。

